% Source: https://artofproblemsolving.com/community/c3h2808581_2018_geometry_bee__preliminary_round
% Stable identity: aops-topic-2808581-q09
Let the arcs (not containing the other vertices) be $AB=\alpha$, $BC=\beta$, $CD=\gamma$, $DA=\delta$, so $\alpha+\beta+\gamma+\delta=360^\circ$, and $\angle ADB=\dfrac\alpha2$.

$AC\perp BD$ gives $\dfrac{\alpha+\gamma}2=90^\circ$, so $\alpha+\gamma=180^\circ$. $\angle DAB=50^\circ$ gives $\beta+\gamma=100^\circ$.

Since $M$ bisects arc $DA$, $\angle MBD=\dfrac\delta4=\angle MCA$, i.e. $\angle YBX=\angle YCX$, so $B,C,X,Y$ are concyclic. Hence the exterior angle $\angle XYM$ of this cyclic quadrilateral at $Y$ equals $\angle XCB=\angle MCB=\dfrac12\left(\alpha+\dfrac\delta2\right)$.
\begin{asy}
usepackage("amsmath");
size(5cm);
defaultpen(fontsize(10pt));
real a=100,bb=20,cc=80,dd=160;pair A=dir(0),B=dir(a),C=dir(a+bb),D=dir(a+bb+cc),M=dir(a+bb+cc+dd/2);
draw(unitcircle);draw(A--C);draw(B--D);draw(C--M);draw(B--M);draw(A--B--C--D--cycle,gray);
label("$A$",A,E);label("$B$",B,N);label("$C$",C,NW);label("$D$",D,SW);label("$M$",M,SE);
\end{asy}

Setting this to $90^\circ$: $\alpha+\dfrac\delta2=180^\circ$. With $\alpha+\gamma=180^\circ$ this gives $\gamma=\dfrac\delta2$; with $\beta+\gamma=100^\circ$ and the total $360^\circ$, $\alpha+\delta=260^\circ$, so $\dfrac\delta2=80^\circ$, $\delta=160^\circ$, $\alpha=100^\circ$. Therefore $\angle ADB=\boxed{50^\circ}$ (confirmed numerically).
